% !PRINTED 401
\weberpage{401}

Ersetzt man $y$ durch eine assoziierte Zahl $\pm\varepsilon^k y$, so geht $\xi_1$
in $\xi_1 + k$ über, worin $k$ eine positive oder negative ganze Zahl
ist, und man kann also $y$ unter den Assoziierten immer auf eine
Weise so wählen, daß

\weq{8}{0 \leqq \xi_1 < 1}

wird. Die so bestimmte Zahl $y$ ist dann die \gesperrt{reduzierte
Zahl}.

Beachtet man noch, daß die Gleichungen (7) sich nicht
ändern, wenn $y$ durch $-y$ ersetzt wird, so folgt:

3. \gesperrt{In einer Schar assoziierter Zahlen gibt es immer
zwei und nur zwei reduzierte Zahlen, die sich nur im
Vorzeichen unterscheiden.}

Aus (7) ergibt sich, wenn $\xi_1$ der Bedingung (8) genügt:

\weq{9}{0 \leqq \log\left|\frac{y}{y'}\right| < \log\frac{\varepsilon}{\varepsilon'},}

und daraus

\weq{10}{1 \leqq \left|\frac{y}{y'}\right| < \frac{\varepsilon}{\varepsilon'},}

und umgekehrt folgen aus (10) wieder die Gleichungen (7) mit
der Bedingung (8).

Da wir überdies \gesperrt{vorausgesetzt haben}, daß $y$ eine \gesperrt{positive}
Norm haben soll, so haben $y$, $y'$ das gleiche Vorzeichen, und dieses
ist positiv, wenn $x_2$ positiv ist. Denn dann ist $y - y' = x_2\sqrt{D}$
positiv, und da $|y| > |y'|$ ist, so müssen $y$ und $y'$ positiv sein.

In dem besonderen Falle $x_2 = 0$ ist $y$ und $y'$ positiv, wenn
wir $x_1$ positiv voraussetzen. Damit ist dann die in 3. noch übriggebliebene
Zweideutigkeit beseitigt.

Da also $y$ und $y'$ positiv sind, so folgt aus (10):

\weq{11}{y'\varepsilon - y\varepsilon' > 0.}

Woraus nach (4) und (6):

\weq{12}{2aUx_1 - (T-bU)x_2 > 0, \qquad x_2 \geqq 0.}

Und diese Bedingungen, die wir die \gesperrt{Isolierungsbedingungen}
nennen, schließen bereits in sich, daß $y'$ und folglich
$N(y)$ positiv ist. Denn nach (3) ist
\[
T > U\sqrt{D},
\]
und demnach folgt aus (12):
\[
U[2ax_1 + (b-\sqrt{D})x_2] > 0.
\]
Damit ist bewiesen:
%CONT>>
